\documentclass[11pt,a4paper]{article}

\usepackage[utf8]{inputenc}
\usepackage[T1]{fontenc}
\usepackage{lmodern}
\usepackage[english]{babel}
\usepackage{geometry}
\usepackage{graphicx}
\usepackage{float}
\usepackage{caption}
\usepackage{amsmath}
\usepackage{amssymb}
\usepackage{booktabs}
\usepackage{hyperref}
\usepackage{subcaption}
\usepackage{xcolor}

\geometry{margin=1in}
\raggedbottom
\emergencystretch=3em

% Figures live in the repo's runs/ directory; this file sits in report/9/,
% so resolve \includegraphics paths relative to the repository root.
\graphicspath{{../../}}

% Graceful fallback for figures produced by runs that may not be pulled yet.
\newcommand{\figorbox}[2]{%
  \IfFileExists{../../#1}{\includegraphics[width=#2\textwidth]{#1}}{%
    \fbox{\parbox[c][3.5cm][c]{#2\textwidth}{\centering\itshape
    \textcolor{gray}{[figure not yet generated --- run the eval and pull it]}}}}}

\title{Empirical Analysis of Transformers on Graph Connectivity \\
       \large Part IX: Beyond the Diameter --- Out-of-Distribution Path Generalization
       and the Role of Endpoints}
\author{Edoardo Paccagnella}
\date{\today}

\begin{document}

\maketitle

% =====================================================================
% DRAFT / PLAN ONLY (2026-07-24). No experiments have been run yet and no
% code beyond what already exists from Reports VI--VIII has been written for
% this report. This file lays out the questions, the recap of what is already
% established, and the two planned experiment threads (A: out-of-distribution
% path generalization, B: the same training regime applied to cycles).
% Every \subsection below under "Planned experiments" is a STUB: it states the
% question and the planned setup, with an explicit "[Data pending --- not yet
% run]" marker in place of results. Do NOT fill in numbers, tables, or figures
% until the corresponding experiment has actually been run on real checkpoints
% and the data has been pulled from HPC (see istruzioni.md, error 2 in the
% error log and the general reading rule at the top of that file).
%
% Design rules carried over unchanged from Reports IV--VIII (istruzioni.md
% §4, especially errors 21, 55--60, 61, 62):
%  - every caption states MODEL + TRAINING SET + TEST SET + METRIC, no analysis;
%  - no code/file/path/flag names in the rendered text (keep them in % comments);
%  - no 'advisor' attributions, no chronological "we re-ran" story;
%  - error bars (std over seeds) on every curve that aggregates seeds;
%  - when a discrete sweep shows something interesting at one point, evaluate
%    the immediately neighbouring points too, not just the next planned one;
%  - every experiment trains on a single, explicitly named distribution, never
%    an opaque random mixture (principle introduced in Report VI, istruzioni.md
%    §13.2).
%  - NEW for this report (explicit instruction from the user, istruzioni.md
%    §16.4): every new training in this report uses canvas size n=46 (not the
%    n=40 used throughout Reports VI--VIII), everything else identical
%    (disjoint-paths distribution where applicable, similarity read-out, same
%    step/sample budget, same RoBERTa-faithful L=2/d512/single-head model).
%    Reuse the existing training/eval code (already parametric in the canvas
%    size), do not reimplement anything from scratch.
% =====================================================================

\begin{abstract}
\noindent Report~VIII showed a standard connectivity transformer, trained on unions of
disjoint paths and evaluated with a similarity read-out, resolving an unbalanced
two-component split far beyond the doubled distance wall $2\cdot 3^L$, and collapsing to
``everything connected'' on a structurally different two-component graph (two cycles) that
the same training distribution never contains. Both results are measured \emph{inside} a
training distribution that already covers every two-segment split of the canvas, so neither
one distinguishes a genuinely general mechanism (the two open \emph{endpoints} of a path)
from plain in-distribution coverage. This report is a plan for two threads that separate the
two readings: (A) a search for a case of path learning that generalises to structure the
model was \emph{never} trained on, and (B) the same systematic training regime applied to
\emph{cycles} (which have no endpoints at all), to test whether endpoints are truly
load-bearing or whether cycles, trained properly, are simply another counterexample to
``the diameter is the difficulty measure.'' No experiment in this report has been run yet;
this document states the questions, the established background, and the planned setups.
\end{abstract}

\section{Introduction}
\label{sec:intro}

\subsection{What Reports I--VIII established}
\label{sec:recap}

We treat the following as established (documented, multi-seed, in Reports~I--VIII); the new
work in this report must stay consistent with them.

\begin{itemize}\setlength\itemsep{0.2em}
  \item \textbf{There is a sharp, architectural distance wall at $3^L$} ($=9$ at $L=2$),
        which moves with depth and appears even on graph families seen $\sim10^8$ times in
        training (Reports~III--IV).
  \item \textbf{The model computes distance-bounded matrix powering, not a bounded
        traversal}, and a soft, learned ``node budget'' on dense structure is a removable
        data prior, not a second capacity (Report~V).
  \item \textbf{Parallel routes and repeated local structure both help, within capacity, in
        data-shaped ways} rather than architecturally for free (Report~VI).
  \item \textbf{A similarity read-out} ($z_{ij}=\mathrm{scale}\cdot\cos(h_i,h_j)+\mathrm{bias}$)
        \textbf{doubles the single-route reach to $2\cdot 3^L$} and makes behaviour far more
        robust to the trunk's imperfect, graded cross-component mixing than the linear
        read-out (Report~IV, confirmed mechanistically in Report~VII).
  \item \textbf{An unbalanced two-component split resolves far beyond the distance wall; a
        balanced split does not, with the same weights} (Report~VI, Thread~B) --- opening the
        trunk (Report~VII) showed why: on a graph with exactly two components, once one is
        small enough to resolve completely, the model gets an extra, graded completion signal
        for the \emph{other} component, built by the same two-layer attention machinery that
        does ordinary beyond-capacity reach, not a separate circuit.
  \item \textbf{Report~VIII collected the full mechanistic battery under the now-standard
        similarity read-out} (behavioural sweep, read-in geometry, attention scores, the exact
        node-to-node contribution, a new causal edge-ablation contribution, layer ablation,
        the three-component falsification test) on the disjoint-paths-trained, $n{=}40$
        checkpoints, and added a stagewise probe locating \emph{where} in the forward pass the
        completion signal is assembled: attention layer~2, not either feed-forward step,
        builds most of the far-reach; MLP~2's distinctive job is sharpening the \emph{cut}.
  \item \textbf{Report~VIII also tested two \emph{cycles} in place of two chains}, on the same
        disjoint-paths-trained checkpoints (which never saw a cycle in training): the model
        collapsed to predicting the entire graph connected at \emph{every} split and every
        seed (predicted-positive rate exactly $1.000$, cut exactly $0.000$) --- read as a
        generic out-of-distribution collapse on an unfamiliar canvas, not a clean test of
        whether path endpoints specifically matter, since the training distribution never
        contained a cycle of any kind.
  \item \textbf{The training distribution behind every two-component result since Report~VI}
        is \texttt{path\_union}-style: an online stream of unions of $1$--$4$ disjoint paths
        that, together, partition all $n$ canvas nodes. Crucially, this stream already
        contains \emph{every} two-segment split of the canvas as a special case (Report~VII,
        \S sec:intro there) --- so any result showing the model resolving a particular
        two-component split is, by construction, an in-distribution result, not evidence of
        generalisation to unseen structure.
\end{itemize}

\subsection{The new question: is the diameter always the right measure, and are path endpoints
load-bearing?}
\label{sec:question}

Report~VIII's two-cycles result and the underlying two-chain completion signal (Report~VI/VII)
together suggest that \emph{diameter alone} is not always the right notion of difficulty: a
model can resolve pairs far beyond $2\cdot 3^L$ on one graph shape and fail completely on a
same-size, same-split graph of a different shape. But every demonstration of ``beyond the
wall'' so far sits \emph{inside} a training distribution that already covers the test
structure. This report asks two questions in sequence, each stated as a plan of experiments to
be run in a later session, not results already in hand.

\begin{enumerate}\setlength\itemsep{0.15em}
  \item \textbf{Is there a case of path learning that is genuinely out-of-distribution?} All
        the evidence so far only shows that, \emph{in distribution}, the model learns beyond
        the doubled wall by using the two open endpoints of a path. Is that a general skill
        that transfers to structure never seen in training (a different diameter cutoff, a
        different number of components, a different internal density), or is it tied to
        having seen exactly this kind of structure during training?
  \item \textbf{Are the two open endpoints of a path truly load-bearing, or is the two-cycles
        failure of Report~VIII simply an artefact of never having trained on a cycle at all?}
        Repeating the same systematic training regime used for paths, but on cycles, and
        reading the identical mechanistic battery, would settle whether cycles are
        architecturally unlearnable in this respect or whether they are simply unfamiliar ---
        in which case cycles become a \emph{second} counterexample to ``the diameter is the
        difficulty measure,'' on equal footing with the path result, rather than evidence that
        endpoints are special.
\end{enumerate}

\subsection{Two candidate readings, stated before running anything}
\label{sec:hypotheses}

As in every previous report, we state the candidate readings up front so the planned
experiments are genuine tests, not a search for evidence of an assumed answer.

\begin{itemize}\setlength\itemsep{0.15em}
  \item \textbf{Reading 1 --- endpoints are load-bearing.} The completion signal is tied
        specifically to the existence of an open, degree-$1$ endpoint; a structure without one
        (a cycle) cannot support it no matter how it is trained, and a path structure stripped
        of its usual training coverage (e.g.\ capped diameter, fewer path counts) loses the
        skill on genuinely unseen splits.
  \item \textbf{Reading 2 --- endpoints are incidental; coverage/familiarity is what matters.}
        Given a training distribution that systematically exposes the model to a given
        two-component shape (paths \emph{or} cycles) across a range of splits, the model
        acquires an analogous completion signal regardless of whether the shape has open
        endpoints; the two-cycles failure of Report~VIII is then simply because cycles were
        never trained on, not because they lack endpoints.
\end{itemize}

Neither reading is assumed; the point of Threads~A and~B below (\S\ref{sec:plan}) is to produce
evidence that could support either one.

\section{Setup}
\label{sec:setup}

\paragraph{Model.} The same base model as Reports~III--VIII: the RoBERTa-faithful standard
transformer, $L=2$, single head, $d_{\mathrm{model}}=512$, normalised-ReLU attention. As in
Report~VIII, we standardise on the \textbf{similarity} read-out
($z_{ij}=\mathrm{scale}\cdot\cos(h_i,h_j)+\mathrm{bias}$) throughout this report, for the same
reason given there: it makes behaviour far more robust to the trunk's imperfect, graded
cross-component mixing than the linear read-out, without removing the underlying mechanistic
signal (Report~VII, \S sec:res-similarity).

\paragraph{Canvas size: $n=46$.} Every new training run in this report uses $n=46$ nodes, not
the $n=40$ used throughout Reports~VI--VIII. Everything else about the training setup is kept
identical to the disjoint-paths condition of Report~VIII (same online distribution family
where applicable, same number of training steps/samples, same architecture, same read-out).
This is an explicit, non-negotiable instruction (\S\ref{sec:plan} restates it per thread where
relevant); the training/evaluation code from Reports~VI--VIII is already parametric in the
canvas size and is reused unchanged, not rewritten.
% NOTE for implementation (not for the rendered report): every existing script that trains or
% evaluates on this project's disjoint-paths / two-chain / cycle families already accepts an
% explicit node-count argument. When Thread A/B experiments are implemented, launch them with
% that argument set to 46 instead of 40; do not hardcode a new value anywhere.

\paragraph{Test graphs.} Each planned experiment below (\S\ref{sec:plan}) specifies its own
test construction (a capped-diameter path union, a two-chain graph with an edge removed to
create a third component, unions of $5$--$7$ disjoint paths, a chain-plus-clique graph, a
two-component cycle graph, and others). No test graphs have been generated yet; the exact
generators to reuse or add are noted in \S\ref{sec:plan} as implementation notes, to be
resolved when each experiment is actually built.

\section{Plan of experiments}
\label{sec:plan}

This report is organised around two threads, in the order the open question in
\S\ref{sec:question} states them. \textbf{No experiment below has been run.} Every subsection
states a question, a planned setup, and stops at a ``\emph{[Data pending --- experiment not yet
implemented]}'' marker; nothing here should be read as a result.

\subsection{Preliminary check --- does canvas size $n=46$ reproduce the effect at all?}
\label{sec:planpre}

Before committing every training run in this report to $n=46$ (\S\ref{sec:setup}), we first
check, on two seeds, that the disjoint-paths/similarity condition at $n=46$ reproduces the same
endpoint-completion signature already established at $n=40$ (Report~VI/VII): a small split
should resolve the long component far past the distance wall, a balanced split should not, with
the same weights. Two independent training runs (seeds $1000$ and $2000$), identical in every
respect to the Report~VIII disjoint-paths condition except the canvas size, followed by the same
split-sweep evaluation used throughout Reports~VI--VIII. \emph{[Data pending --- training
launched by the user, not yet returned; code prepared and reused unchanged from
Reports~VI--VIII.]}

\subsection{Thread A --- out-of-distribution generalisation for path-trained models}
\label{sec:planA}

Thread~A is the more important of the two (explicitly, in the request behind this report): find
at least one case where a path-trained model's completion signal is tested on structure it
provably never saw, rather than on a special case already covered by its training stream.

\paragraph{A.1 --- Train with a capped diameter, test beyond it.} Train on a union-of-paths
distribution restricted to a maximum diameter (e.g.\ two-chain splits that never exceed
diameter $18$, the doubled wall $2\cdot 3^L$), then evaluate on two-chain splits whose diameter
\emph{exceeds} that cap --- structure the training stream never produced by construction. This
is the most direct test of whether the completion signal generalises past what the model has
been shown, or whether it is bounded exactly by what was covered in training.
\emph{[Data pending --- not yet implemented.]}

\paragraph{A.2 --- Less variety in the training distribution.} The current disjoint-paths
stream (unions of $1$--$4$ path components) may be more varied than necessary for the model to
acquire the endpoint-completion skill. Train instead on a narrower distribution (e.g.\ only
two-chain splits, never $3$ or $4$ components) and test whether the completion signal still
appears. If it does, the phenomenon does not depend on the richness of the training
distribution (a more general, more interesting result); if it does not, the phenomenon is
specific to the richer stream (a narrower but still reportable result).
\emph{[Data pending --- not yet implemented.]}

\paragraph{A.3 --- Genuinely unseen three-component test.} Train \emph{only} on two-chain
splits (never three or more components), then evaluate on a graph built by removing one
internal edge from a two-chain test graph, producing three components that were never a
possibility under this training distribution (unlike the three-component test already run in
Reports~VII--VIII, whose training distribution could already produce three or four components).
\emph{[Data pending --- not yet implemented.]}

\paragraph{A.4 --- Existing checkpoints on 5/6/7-component unions.} Using the already-trained
path\_union checkpoints (which only ever see $1$--$4$ components in training), evaluate on
unions of $5$, $6$, and $7$ disjoint paths --- component counts never produced during training
--- to see whether the endpoint-completion skill generalises to a number of components beyond
what was covered. One long component plus $K{-}1$ short ones, with the short-component size
chosen so the long component's internal diameter exceeds $18$ (the doubled wall), so a failure
cannot be attributed to the ordinary distance wall alone. This experiment is eval-only on
existing $n{=}40$ checkpoints and is the one exception to the $n{=}46$ rule of
\S\ref{sec:setup} (no retraining involved). \emph{[Data pending --- eval code implemented and
smoke-tested, launched by the user, not yet returned.]}

\paragraph{A.5 --- Two-chain-trained checkpoints on targeted structures.} Train (four seeds, as
throughout this project) on a range of two-chain splits only, confirm the endpoint-completion
signal reappears, then evaluate the resulting checkpoints on structures chosen to isolate what
actually matters: (i) a graph made of one chain and one clique (two components, only one of
them with a genuine open endpoint); (ii) a two-component graph shaped like a two-chain split but
where the \emph{internal} (non-endpoint) nodes of each component carry additional edges among
themselves, so the two components stay separate and still have clearly identifiable endpoints
but are internally denser. The goal is to separate ``needs the exact path shape'' from ``needs
two components with recognisable endpoints.''
\emph{[Data pending --- not yet implemented.]}

\subsection{Thread B --- are endpoints truly necessary? The analogous experiment on cycles}
\label{sec:planB}

Thread~B asks the mirror question directly: instead of treating the two-cycles collapse of
Report~VIII as evidence that endpoints are necessary, train the model on cycles the same
systematic way paths were trained, and see whether an analogous completion signal appears.

\paragraph{B.1 --- The same training regime as paths, applied to cycles.} Repeat, with cycles
in place of paths, the same kind of systematic training distribution used for the path results
of Reports~VI--VIII (an explicit, named distribution over cycle unions/splits, not an opaque
mixture --- consistent with the data principle introduced in Report~VI, \S13.2 of the project
instructions), sweeping the split the same way, and evaluate \emph{in-distribution} (the same
kind of evaluation the path results received, as opposed to the purely out-of-distribution
two-cycles test of Report~VIII, where the training stream never contained any cycle at all).
The aim is to see whether, and how, the model learns to handle cycles when they are not a
foreign structure.
\emph{[Data pending --- not yet implemented.]}

\paragraph{B.2 --- The full Report~VIII mechanistic battery, repeated on cycle-trained
checkpoints.} On the resulting cycle-trained checkpoints, repeat the entire battery assembled in
Report~VIII: the general behavioural results, the read-in geometry analysis, the attention
scores, the exact node-to-node contribution, the layer ablation, a three-cycles falsification
test (the structural analogue of the three-component split test), a two-cycles split sweep, and
the layerwise cosine geometry (stagewise probe). The purpose is to determine, before concluding
that path endpoints are important or that cycles are intrinsically harder than other graph
shapes, whether cycles --- trained on properly --- can \emph{also} serve as a counterexample to
``the diameter is the difficulty measure,'' on the same footing as the path result.
\emph{[Data pending --- not yet implemented.]}

\section{Status}
\label{sec:status}

Most of \S\ref{sec:plan} is still unimplemented. Two exceptions, both code-complete and
smoke-tested but with no real data back yet: the preliminary $n{=}46$ check
(\S\ref{sec:planpre}, two training runs) and Thread~A.4 (\S\ref{sec:planA}, an eval-only test on
existing $n{=}40$ checkpoints). Everything else in this document exists to fix the questions,
the established background, and the exact planned setups before any further code is written, so
that the next session can proceed thread by thread --- one experiment at a time, each
smoke-tested on a toy checkpoint before touching real weights, with every number in a future
revision of this report pulled from real data on real checkpoints, never written from memory or
from partial output.

\end{document}
